%% ============================================================================
%% MASTER CV (compact style) - [YOUR_NAME]
%%
%% This file is the FULL master content bank (applications/master_cv.md) rendered
%% in the compact LaTeX style. It is the LaTeX skeleton AND the verbatim
%% content source for tailored CVs.
%%
%% HOW TO BUILD A TAILORED CV (verbatim selection):
%%   1. Copy this file to applications/<NN>_<company>_<role>/CV_<CVNameSlug>_<company>_<role>.tex
%%   2. SELECT + REORDER: keep the blocks relevant to the role, delete the rest,
%%      reorder within each section. Content is verbatim-FIRST: projects,
%%      education lines, and headers stay character-for-character.
%%      Experience bullets may be LIGHTLY rephrased when it genuinely improves
%%      role fit (tightening, aligning to the posting's terminology), but facts,
%%      metrics, and scope stay exactly the master's - no new claims, no
%%      escalated numbers, no merged achievements. Verbatim stays the default.
%%   2a. TWO surfaces are composed per role, never copied: the About Me
%%      paragraph (fully generative) and the CORE COMPETENCIES section (5-7
%%      \cvcomp bullets written for THIS posting). Core Competencies REPLACES
%%      the Skills bank below - a tailored CV prints Core Competencies and
%%      deletes \section{Skills} entirely. Every tool, method, and metric named
%%      in a competency bullet must trace to this master or to the candidate
%%      profile; the composing, labelling, and ordering are what is free.
%%   2b. EVERY project entry carries a mandatory one-line "Tech stack:" subtitle
%%      (\cvproject arg 2). Stacks trace to this master's own stack line, trimmed
%%      to the most role-relevant items so the line never wraps. Keep it under
%%      ~95 characters; drop items rather than let it run to two lines.
%%      Never repeat stack items in the project title or description.
%%   2c. SELECTED PROJECTS closes with a mandatory one-line \cvmoreprojects
%%      "Also built:" catch-all naming the master projects that were cut, so the
%%      selection reads as a selection rather than the whole inventory.
%%   3. Keep the \newpage before CORE COMPETENCIES: tailored CVs are exactly
%%      2 pages - page 1 = About Me / Education / Work Experience / Languages,
%%      page 2 = Core Competencies / Selected Projects / Other Relevant
%%      Information. There is no standalone References section: its line is
%%      the mandatory first bullet of Other Relevant Information.
%%   4. The About Me paragraph must END with the portfolio sentence:
%%      Get to know me better \href{<your portfolio URL>}{here}!
%%      Delete the sentence if you have no portfolio.
%%   5. Self-description in the About Me opening should default to
%%      [YOUR_PREFERRED_SELF_DESCRIPTION] - pick a variant that fits the posting
%%      without overstating. Respect any framing-suppression rule configured in
%%      .claude/skills/job-application-assistant/03-writing-style.md.
%%
%% This master compiles to several pages by design (it is the full bank).
%% Only tailored CVs have the 2-page rule.
%%
%% SETUP: /setup replaces the bracketed placeholders in this file with your
%% actual details. Placeholders in LIVE CODE use spaces, not underscores, so the
%% file compiles before setup runs ([YOUR_NAME] with an underscore appears only
%% in comments and metadata, where LaTeX treats it as safe text).
%%
%% Compile with lualatex (fontspec + system Carlito font):
%%   cd applications && lualatex -interaction=nonstopmode -output-directory=<NN>_<company>_<role> <NN>_<company>_<role>/CV_<CVNameSlug>_<company>_<role>.tex
%% ============================================================================

\documentclass[10pt,a4paper]{article}

\usepackage[a4paper,left=2.0cm,right=2.0cm,top=1.6cm,bottom=1.7cm]{geometry}
\usepackage{fontspec}
\setmainfont{Carlito}   % metric-compatible Calibri clone; clean ATS text layer
\usepackage{xcolor}
\definecolor{cvblue}{HTML}{2E74B5}   % Word heading blue
\definecolor{cvdark}{HTML}{262626}   % near-black for the contact line
\usepackage{enumitem}
\usepackage{titlesec}
\usepackage{needspace}
\usepackage{fontawesome5}   % LinkedIn/GitHub/globe brand glyphs for the contact line
\usepackage{hyperref}
\hypersetup{
    colorlinks=true,
    urlcolor=cvblue,
    linkcolor=cvblue,
    pdftitle={[YOUR_NAME] - CV},
    pdfauthor={[YOUR_NAME]},
}

\pagestyle{empty}
\setlength{\parindent}{0pt}
\setlength{\parskip}{5pt}
\emergencystretch=1em
\linespread{1.04}
\frenchspacing

% Section heading: ALL-CAPS bold blue with a thin full-width rule hugging the text
\titleformat{\section}
  {\color{cvblue}\bfseries\fontsize{11}{13}\selectfont}
  {}{0pt}{\MakeUppercase}
  [\vspace{-10pt}{\color{cvblue!55}\rule{\textwidth}{0.5pt}}]
\titlespacing*{\section}{0pt}{8pt}{0pt}

% One-line work header: \cvjob{Job Title, Company}{Mon YYYY - Mon YYYY}
\newcommand{\cvjob}[2]{\needspace{4\baselineskip}\vspace{3pt}{\bfseries #1 \textbar{} #2\par}\vspace{-2pt}}
% Optional italic context line directly beneath a work header
\newcommand{\cvcontext}[1]{{\itshape\small\hspace{0.3em}#1\par}\vspace{-3pt}}
% One-line education entry: \cvedu{Degree}{Institution}{Mon YYYY - Mon YYYY}
\newcommand{\cvedu}[3]{{\bfseries #1} \textbar{} #2 \textbar{} #3\par}
% Optional italic thesis/final-project line directly beneath an education entry
\newcommand{\cvedusub}[1]{\vspace{-4pt}{\itshape\small\hspace{0.3em}#1\par}}
% Bare blue separator rule (used where a section has no heading, e.g. above About Me)
\newcommand{\cvrule}{\par\vspace{2pt}{\color{cvblue!55}\rule{\textwidth}{0.5pt}}\par}
% Skills row: \cvskill{Category}{item · item · item}
% MASTER-ONLY in tailored CVs: these rows are the evidence bank Core Competencies
% is composed from, not a section a tailored CV prints.
\newcommand{\cvskill}[2]{\par{\bfseries #1:} #2\par}
% Core-competency bullet: \cvcomp{Label}{Concrete items, then what it buys the role.}
% Used INSIDE an itemize. Composed per role - 5-7 bullets, label carries the
% posting's own core term wherever that is truthful.
\newcommand{\cvcomp}[2]{\item {\bfseries #1:} #2}
% Catch-all closing line for Selected Projects: \cvmoreprojects{short list. Portfolio link.}
% ONE rendered line (two at the absolute most); names only master projects that
% were not selected above.
\newcommand{\cvmoreprojects}[1]{\needspace{2\baselineskip}\vspace{2pt}{\bfseries Also built:} #1\par}
% Project entry: \cvproject{Title}{Tech stack}{Description}
% The tech-stack line is MANDATORY and must stay on ONE line (keep it under ~95
% characters including the "Tech stack: " label). Separate items with " · ".
% Because the stack line carries the tooling, the title and description must NOT
% repeat it - no "(Python + Ollama)" in the title, no "built with X, Y, Z" prose.
\newcommand{\cvproject}[3]{\needspace{4\baselineskip}\vspace{2pt}{\bfseries #1}\par\vspace{-4pt}{\itshape\small Tech stack: #2\par}\vspace{-4pt}#3\par}

% Compact bullet list
\setlist[itemize]{label=\textbullet,leftmargin=1.15em,itemsep=1.5pt,parsep=0pt,topsep=2pt,partopsep=0pt}

\begin{document}

% ============================================================
%   HEADER  (fixed contact line; no headline)
% ============================================================

{\fontsize{20}{24}\selectfont\bfseries\color{cvblue}[Your Full Name]\par}
\vspace{1pt}
{\bfseries\color{cvdark}[Nationality] \textbar{} [+00 000 000 000] \textbar{} your.email@example.com \textbar{} \href{https://your-portfolio.example.com/}{\faGlobe\;your-portfolio.example.com} \textbar{} \href{https://www.linkedin.com/in/your-profile}{\faLinkedin\;your-profile} \textbar{} \href{https://github.com/your-username}{\faGithub\;your-username}\par}

% ============================================================
%   ABOUT ME  (no heading - separator rule only; the only fully
%   generative content in a tailored CV)
% ============================================================

\cvrule

[2-4 sentence summary]: who you are professionally, the throughline of your experience, and your strongest evidence of it. Lead with the most role-relevant point and close with the broader profile as context. End with the portfolio sentence: Get to know me better \href{https://your-portfolio.example.com}{here}!

% ============================================================
%   EDUCATION  (one bold line each + optional thesis/final-project
%   sub-line; full project descriptions still live in Projects)
% ============================================================

\section{Education}

\cvedu{[Degree]}{[Institution]}{Sep 20XX - Jul 20XX}
\cvedusub{Thesis: [Thesis or final-project title] ([grade])}
\cvedu{[Degree]}{[Institution]}{Sep 20XX - Jul 20XX}
\cvedusub{Final Project: [Final-project title] ([grade])}
%% Add or delete \cvedu/\cvedusub pairs as needed; early or optional entries can
%% stay off tailored CVs unless the posting calls for them.

% ============================================================
%   WORK EXPERIENCE  (one block per role, most recent first.
%   Tailored CVs select 5-6 bullets from the most recent role,
%   3-4 from the second, 2-3 from earlier roles.)
% ============================================================

\section{Work Experience}

\cvjob{[Job Title], [Company]}{Mon 20XX - Mon 20XX}
\cvcontext{[One-line company/product context - what the company does and at what scale.]}
\begin{itemize}
\item [Achievement or responsibility bullet, quantified wherever possible.]
\item [Achievement or responsibility bullet.]
\item [Achievement or responsibility bullet.]
\end{itemize}

\cvjob{[Job Title], [Company]}{Mon 20XX - Mon 20XX}
\cvcontext{[One-line company/product context.]}
\begin{itemize}
\item [Achievement or responsibility bullet.]
\item [Achievement or responsibility bullet.]
\end{itemize}

\cvjob{[Earlier Job Title], [Company]}{Mon 20XX - Mon 20XX}
\begin{itemize}
\item [Achievement or responsibility bullet.]
\item [Achievement or responsibility bullet.]
\end{itemize}

% ============================================================
%   LANGUAGES  (one line, closes page 1 in tailored CVs)
% ============================================================

\section{Languages}

[Language]: [Level] · [Language]: [Level] · [Language]: [Level]

% ============================================================
%   PAGE BREAK - tailored CVs keep this so page 2 = Core Competencies onward
% ============================================================
\newpage

% ============================================================
%   CORE COMPETENCIES  (COMPOSED PER ROLE - the bullets below are a
%   WORKED EXAMPLE, NOT content to copy. A tailored CV rewrites all of
%   them for its own posting: 5-7 bullets, bold label carrying the
%   posting's own core term where truthful, body listing concrete
%   items drawn from the SKILLS bank below and closing with what the
%   competency buys THIS role. Every named tool, method, and metric
%   must trace to this master or to the candidate profile.)
% ============================================================

\section{Core Competencies}

\begin{itemize}
\cvcomp{[Competency label]}{[Concrete tools, methods, and metrics drawn from the Skills bank below, then one clause on what this buys the role.]}
\cvcomp{[Competency label]}{[Concrete items and value clause.]}
\cvcomp{[Competency label]}{[Concrete items and value clause.]}
\cvcomp{[Competency label]}{[Concrete items and value clause.]}
\cvcomp{[Competency label]}{[Concrete items and value clause.]}
\cvcomp{[Competency label]}{[Concrete items and value clause.]}
\end{itemize}

% ============================================================
%   SKILLS  (SOURCE POOL - the evidence bank Core Competencies is
%   composed from. A tailored CV DELETES this whole section; it is
%   never printed alongside Core Competencies.)
% ============================================================

\section{Skills}

\cvskill{[Skill Category 1]}{[item · item · item · item] }
\cvskill{[Skill Category 2]}{[item · item · item · item]}
\cvskill{[Skill Category 3]}{[item · item · item · item]}
\cvskill{[Skill Category 4]}{[item · item · item · item]}
\cvskill{[Skill Category 5]}{[item · item · item · item]}
\cvskill{[Tools \& Platforms]}{[tool · tool · tool · tool]}

% ============================================================
%   SELECTED PROJECTS  (always include; select 4-6 per role in
%   tailored CVs, then close with the mandatory \cvmoreprojects
%   "Also built:" line naming what was cut)
% ============================================================

\section{Selected Projects}

\cvproject{[Project Title 1]}{[Tech stack item · item · item]}{[What the project is and what came of it - outcome, scale, or validation.]}

\cvproject{[Project Title 2]}{[Tech stack item · item · item]}{[What the project is and what came of it.]}

\cvproject{[Project Title 3]}{[Tech stack item · item · item]}{[What the project is and what came of it.]}

\cvproject{[Project Title 4]}{[Tech stack item · item · item]}{[What the project is and what came of it.]}

\cvproject{[Project Title 5 (optional)]}{[Tech stack item · item · item]}{[What the project is and what came of it.]}

% --- Mandatory closing catch-all. In a tailored CV this line names the master
% --- projects that did NOT make the selection above, so the section reads as a
% --- selection. One rendered line; titles compressed to short descriptors;
% --- never a project absent from this master.
\cvmoreprojects{[a cut project, a second cut project, a third cut project]. More at \href{https://your-portfolio.example.com}{your-portfolio.example.com}.}

% ============================================================
%   OTHER RELEVANT INFORMATION  (closes the CV; select 2-4 of the
%   optional bullets in tailored CVs. The references bullet is
%   MANDATORY and always leads the section. It never prints
%   referee names or contact details.)
% ============================================================

\section{Other Relevant Information}

\begin{itemize}
\item [YOUR REFERENCES LINE - e.g. Recommendation letter from <reference name and title>, available on request; additional references on request.]
\item [Optional bullet: languages, certifications, volunteering, international exposure.]
\item [Optional bullet: awards, publications, community involvement.]
\item [Optional bullet: personal interests that genuinely add signal for this role.]
\end{itemize}

\end{document}
